% Primary-text treatment of planned 164.3; the distinct Forster2011 facet remains separate.
\chapter{Whewell and Mill on Scientific Induction}
\label{sec:p15-whewell-mill}

Does discovering a form which connects observations already constitute an
induction, or must induction carry an inference beyond what was observed?
William Whewell and John Stuart Mill disagree about this question while
agreeing that scientific inquiry requires both inventive thought and exact
comparison with facts. Their disagreement is not captured by assigning
hypotheses to Whewell and observations to Mill.

The primary presentations used here are Whewell's \emph{Novum Organon
Renovatum}, the second part of the third edition of his \emph{Philosophy
of the Inductive Sciences}, published in 1858, and the 1882 Harper
printing of the eighth edition of Mill's \emph{System of Logic}.
These are distinct dated presentations, not the unmodified first editions.
The selected Whewell chapters concern facts, colligation and the
characteristics of induction; the selected Mill passages concern
induction, causation, experimental methods and plurality of causes
\parencite[Book~II, chs.~III--V, pp.~50--96]{Whewell1858NovumPG}
\parencite[Book~III, chs.~I--II, V, VIII and X]{Mill1882SystemLogicPG}.
Their methodological and historical arguments are not experimental findings
about how every person reasons, and neither author's proposed test is
automatically a modern theorem of statistical identification.

\section{Facts, conceptions and colligation}
\label{sec:p15-wm-colligation}

Whewell rejects the proposal to make facts certain simply by removing every
mental contribution. Recognizing an inscription, describing a tree as moved
by a force, identifying a shape and counting its sides already involve
interpretation. His windmill, ambiguous-diagram and apparent-lunar-size
examples expose errors in what initially looks like pure perception.
His response is not to abandon observation: distinguish the Ideas being
applied, make their application precise, and decompose a complex appearance
into definite elementary facts
\parencite[II.III.1--3, pp.~50--54]{Whewell1858NovumPG}.

In the subsequent articles he requires intellectual rather than portentous
or emotional descriptions, precise spatial and temporal relations where
appropriate, and equally definite conceptions outside geometry.
Force, musical concord and resemblance are not reducible to merely counting
objects. Decomposition also introduces usable terms and classifications.
It is necessary but not sufficient: Whewell's contemporary weather records
and the counting of crystal faces illustrate numerical precision without
the conception which gives the relevant connection. The later recognition
that angles between crystal faces matter is his example of a further step.
These are his historical assessments, not a claim that modern meteorology
has no laws
\parencite[II.III.4--11, pp.~54--58]{Whewell1858NovumPG}.

\begin{hbdefinition}{def:p15-wm-colligation}{Whewell's colligation and induction}
In Whewell's presentation, \emph{colligation of facts} is the intellectual
operation which connects phenomena by an appropriate conception, yielding
a precise general proposition. \emph{Induction} names both a true
colligation by an exact and appropriate conception and its resulting
proposition. It is not merely the juxtaposition or verbal sum of records:
the connection includes a conception supplied by the mind.
\end{hbdefinition}

For example, planetary periods and distances are separate quantities until
squaring the periods, cubing the distances and conceiving their
proportionality connects them in Kepler's third law. Whewell then treats
a central accelerating force as a further conception connecting those
relations. This is his account of the inferential steps; the Newtonian
dynamical derivations are not proved merely by naming the central force
\parencite[II.IV.1, pp.~59--60; II.V.1--3, pp.~71--76]{Whewell1858NovumPG}.

The conception may be a new use of an existing idea, not a novel word.
Whewell contrasts explanations using pressure on a lever, weight of air,
refraction and interference with the inadequate earlier conceptions.
Changing a geometrical arrangement within an old hypothesis need not find
the mechanical or optical idea the problem requires.
He describes Kepler's nineteen unsuccessful hypotheses for Mars as
evidence of inventive trial, not as a fixed nineteen-step algorithm.
Earlier investigators' failures may also leave later investigators a
better-defined problem
\parencite[II.IV.6--9, pp.~63--66]{Whewell1858NovumPG}.

Both invention and constraint matter. Whewell says that appropriate
conjectures cannot be guaranteed by rules; he nevertheless requires that
their consequences be steadily worked out and compared with the facts.
He uses calculated lunar places and musical-string relations as
testable consequences. A conception can initially arrive clothed in
dispensable machinery: particles with poles for polarization, a fluid
image of heat, gravitational ether or atoms. The success of a relation
must therefore be distinguished from the literal truth of every image
which first made the relation intelligible
\parencite[II.IV.10--12, pp.~66--69]{Whewell1858NovumPG}.

\section{The Kepler dispute: finding a predicate and extending its scope}
\label{sec:p15-wm-kepler}

Whewell distinguishes discovering the predicate \emph{elliptical motion}
from extending a proposition about particular planets to all planets.
His discussion of the planet-ellipse syllogism asks how that predicate
was first invented and verified. It does not establish a deductively
valid universal conclusion from observations of three planets.
He explicitly separates the inventive step in discovery from verification
in its proof. A definition and a proposition may emerge together:
the definition expresses the conception, the proposition its factual
application. Once accepted, the conception can become so familiar that
its earlier intellectual contribution is no longer noticed
\parencite[II.V.3--5, pp.~74--78]{Whewell1858NovumPG}.

Mill's criterion is different. He requires a transition from cases known
to cases not known. A complete enumeration, or a shorter sentence which
merely collects already established singular statements, does not make
that transition. His examples contrast all the already examined planets
with all planets, including unexamined cases, and distinguish a
mathematical demonstration applicable to every case from conjecturing
a law by checking a few terms
\parencite[III.I.1--2; III.II.1--2]{Mill1882SystemLogicPG}.

\begin{hbdefinition}{def:p15-wm-mill-induction}{Mill's inductive extension}
For Mill, induction infers that what is known in particular cases holds
in other cases resembling them in assignable respects, including a
general proposition's unobserved instances.
Colligation can organize known facts without that extension.
Induction also colligates, but not every colligation is an induction.
\end{hbdefinition}

Mill's navigator sails around a coast and calls the resulting observed
whole an island. In his analysis the word sums up the successive
observations; it does not infer a further coast which was not observed.
He gives Kepler's ellipse the same descriptive standing, while
recognizing that selecting it was much harder. The extension to intervening
and future planetary positions is the genuinely inductive part.
Mill argues that the already established periodic recurrence supplied that
extension, whereas the new ellipse changed the predicate used to express
the path
\parencite[III.II.3 and 5]{Mill1882SystemLogicPG}.

This is Mill's argument about the example, not an independently
reconstructed history of Kepler's calculations. Apparent positions,
the geometric representation of a planetary path and the background
assumptions used to infer that path must not be treated as identical
observational objects.
Mill nevertheless accepts that a conception is required and that its
selection proceeds by tentative, skilful guesses. He rejects the inference
that conceptual involvement alone makes the operation inductive.
His stated objection is to identifying invention with proof, not to
inventive work itself
\parencite[III.II.3--5]{Mill1882SystemLogicPG}.

The differences remain substantive. For Whewell, an appropriate new
connection is part of how general knowledge is constituted. For Mill,
the key logical question is what warrants extending the claim beyond
the known cases. Whewell's selected chapters distinguish invention and
verification; Mill's accusation that the rival account overlooks proof
must not erase that distinction.
Nor does Mill deny the usefulness of colligation: he calls it a necessary
preparation for induction concerning phenomena observed only piecemeal
\parencite[II.V.3]{Whewell1858NovumPG}
\parencite[III.II.4--5]{Mill1882SystemLogicPG}.

Mill's note~105 adds another disagreement. Whewell, quoted there from
his separate \emph{Philosophy of Discovery}, allows parts of older
explanations to survive when their laws are brought into agreement with
a later theory. Mill replies that incompatible assertions about a
physical mechanism cannot both be literally true merely because they
represent the same motions. The quotation is evidence for Mill's
presentation of that reply; it is not a complete reading of the
separate Whewell book
\parencite[III.II, note~105]{Mill1882SystemLogicPG}.

\section{Verification, prediction and consilience}
\label{sec:p15-wm-consilience}

Whewell's account does not infer literal truth from useful fit alone.
Epicycles can organize motions and predict positions while giving a false
picture of the heavens. The older vacuum-avoidance account can organize
some effects without being the final causal account.
He allows false or overburdened hypotheses to contribute useful
connections and to be simplified later
\parencite[II.V.6--10, pp.~78--87]{Whewell1858NovumPG}.

He distinguishes three tasks: representing the facts used in construction;
predicting unobserved cases of the same kind; and accounting for a
different class of facts not contemplated in constructing the theory.
The third is especially important to his stronger test.

\begin{hbdefinition}{def:p15-wm-consilience}{Whewell's consilience}
\emph{Consilience of inductions} occurs when inductions obtained from
different classes of facts coincide within the same theory.
In Whewell's crucial stronger comparison, a hypothesis constructed for
one class supplies the rule and reason of another, previously
uncontemplated class without an adjustment made for that purpose.
Repeatedly adapting a fit to the same class is not that comparison.
\end{hbdefinition}

The worked connections are integral to his argument
\parencite[II.V.10--11, pp.~85--90]{Whewell1858NovumPG}:

\begin{center}
\begin{tabular}{@{}p{.42\linewidth}p{.50\linewidth}@{}}
\toprule
Earlier inferential setting & Further connection reported by Whewell\\
\midrule
Inverse-square attraction and Kepler's third law &
The first and second Kepler laws fall under the same dynamical account.\\
Gravitation inferred through lunar and planetary perturbations &
Precession is accounted for despite its apparently remote character.\\
Interference developed for fringes &
Thin-plate colours are explained without Newton's separate fits apparatus.\\
Spheroidal undulations for double refraction &
The directions of refracted rays connect with their polarization planes.\\
Optical theory extended to polarization &
Brewster's angle, polarized Newton rings and dipolarized-light phenomena
provide further connections.\\
Wavelength estimated from fringes &
The result nearly agrees with the independent thin-plate-colour estimate.\\
\bottomrule
\end{tabular}
\end{center}

These are historical arguments reported by Whewell, not newly reproduced
optical experiments or proofs of every physical assertion.
He makes a strong truth claim for consilience, but his statement that he
knows no counterexample carries the qualification of his historical
knowledge. No theorem here converts that survey claim into a universal
guarantee of infallibility.

Progressive simplification supplies the related comparison.
A successful theory extends its scope while additional suppositions
resolve into existing ones; a failing theory accumulates disconnected
machinery. Whewell contrasts added epicycles with ellipses; Cartesian
vortices with universal gravitation; phlogiston supplemented by an
unusual weight assumption with the oxygen account; and successively
added emission machinery with the wave account's connected extensions.
His criticism of Laplace's gas theory concerns assumptions added for
latent heat which do not simplify the earlier calculations.
This is simplicity under expanding explanatory scope, not merely
preferring the initially shortest formula
\parencite[II.V.12--13, pp.~90--96]{Whewell1858NovumPG}.

\section{Mill's causal question and experimental methods}
\label{sec:p15-wm-methods}

Mill limits this inquiry to physical causes, not ultimate efficient
or ontological causes. The complete cause includes the relevant
concurrence of positive conditions and absence of preventing conditions.
Ordinary speech may single out the last event or the condition needing
explanation; that usage does not make it the entire scientific cause.
His discussion of agents and patients likewise denies that the latter
contribute nothing merely because the sentence calls them passive
\parencite[III.V.2--4, including notes~114--115]{Mill1882SystemLogicPG}.

His further requirement is unconditional succession under the relevant
constitution of nature, not simply an unbroken succession in past
observations. Night preceding day does not suffice: illumination depends
on other conditions. His note~116 distinguishes ultimate laws from
contingent collocations; note~117 qualifies the convenient straight-line
account by refraction. Capacities acquired by an object also require
the later conditions in which they issue in an effect; preparing a
capacity and producing its eventual effect are different sequences
\parencite[III.V.5--6 and notes~116--119]{Mill1882SystemLogicPG}.

In the following source schemata, capital letters name antecedent
circumstances and lower-case letters consequent phenomena. Juxtaposition
lists circumstances or phenomena; it is not multiplication.
Mill's chapter is entitled four methods, but has five canons because
the joint method is a modification joining agreement and difference.

\begin{hbmethod}{def:p15-wm-agreement}{Agreement}
Compare positive instances and eliminate circumstances which are absent
in some instance despite the phenomenon's presence.
The source scheme compares \(ABC\) followed by \(abc\) with \(ADE\)
followed by \(ade\); \(A\) and \(a\) alone remain common.
The first canon identifies the uniquely common circumstance as cause
or effect, under the simplified causal setting.
The two inquiry directions are distinct: finding effects of a known
agent, and finding a cause of an observed effect.
\end{hbmethod}

Mill illustrates the former by oil and alkali producing soap under
varied circumstances, and the latter by solidification from a liquid
state preceding crystallization. He stresses artificial production
as a further test and reports quartz and artificial-marble cases.
The enumeration of antecedents can still be incomplete: an apparently
unique common antecedent may conceal another.
Chapter~X adds the stronger failure from alternative causes
\parencite[III.VIII.1 and 3; III.X.1--2]{Mill1882SystemLogicPG}.

\begin{hbmethod}{def:p15-wm-difference}{Difference}
Compare a positive and a negative instance alike in every materially
relevant respect except one. In the scheme, \(ABC\) produces \(abc\),
whereas \(BC\) produces \(bc\).
The second canon makes the differing circumstance the effect, cause,
or indispensable part of the cause of the phenomenon.
It does not conclude that \(A\) is sufficient by itself independently
of \(B\), \(C\) or other fixed conditions.
\end{hbmethod}

Mill gives this method greater direct causal force and explains why
deliberate experiments can sometimes secure the required comparison
more readily than spontaneous natural cases. Circumstances already
known to be irrelevant need not be identical. The manner of introducing
a change can itself introduce a cause of the result, so it must be
separately checked. A short before--after interval is not, merely by
being short, a proof that no relevant circumstance changed
\parencite[III.VIII.2--3]{Mill1882SystemLogicPG}.

\begin{hbmethod}{def:p15-wm-joint}{Joint agreement and difference}
Compare varied cases with the phenomenon and varied cases without it.
The third canon asks for a uniquely common present circumstance in
the positive cases and its uniquely common absence in the negative cases.
This is Mill's indirect method of difference.
\end{hbmethod}

It is not automatically the direct paired experiment.
Mill explains the difficulty of knowing that no unnoticed common
antecedent remains. His double-refraction discussion connects the
phenomenon with crystalline structure without asserting that every
crystal doubly refracts. Chapter~X explains how the negative comparison
can address plurality of causes, but also why exhausting the negative
field is difficult: the nontransparent substances are not a conveniently
bounded class merely because transparent ones are easier to compare
\parencite[III.VIII.4; III.X.2]{Mill1882SystemLogicPG}.

\begin{hbmethod}{def:p15-wm-residues}{Residues}
From a compound phenomenon subtract portions already attributed to
known antecedents by prior inductions. In the source scheme \(ABC\)
produces \(abc\); if the effects of \(A\) and \(B\) are already known
as \(a\) and \(b\), respectively, the residual \(c\) is referred to
the remaining antecedent \(C\).
This fourth canon uses a deduced negative comparison, not necessarily
a newly observed instance without \(C\).
\end{hbmethod}

The previous effect attributions and their applicability in combination
must be warranted. \(C\) must exhaust the unexplained antecedents;
otherwise the residue identifies an outstanding problem rather than
a uniquely established cause. Mill calls for separate artificial
production or a deductive account from known laws to complete the evidence.
The method is valuable precisely because a small unexplained residue
may expose an inconspicuous agent which was not sought initially
\parencite[III.VIII.5]{Mill1882SystemLogicPG}.

\begin{hbmethod}{def:p15-wm-variation}{Concomitant variations}
Where a permanent agent cannot be removed, compare modifications of it
and the associated changes. The fifth canon permits a cause, an effect,
or a connection through another causal fact; it does not orient the
causal relation by covariance alone.
Hold the other relevant antecedents fixed or otherwise warrant the
comparison which attributes the change.
\end{hbmethod}

Mill contrasts moving a pendulum away from a mountain with the
impossibility of removing the earth, and discusses heat and expansion,
lunar position and tides, and the direction of terrestrial falling.
Correspondence need not mean both quantities move in the same direction.
He explicitly allows two varying phenomena to be effects of a common
cause. Controlled production, where possible, is an additional test
of direction, not a consequence of observing the variations.

Even numerical extrapolation requires qualifications. The absolute
quantities must be known, not just their changes; unobserved counteracting
agents may appear; and different laws can give nearly identical values
over a narrow observed range. Mill contrasts heat-distance speculation
with the diminishing-retardation argument about motion, and retains
Herschel's warning about extrapolated steam and fluid-resistance formulae
\parencite[III.VIII.6--7 and note~136]{Mill1882SystemLogicPG}.

\section{Two explicit boundaries on causal and predictive conclusions}
\label{sec:p15-wm-boundaries}

Mill's own plurality-of-causes argument defeats a reading of agreement
as an unconditional inference from the table alone.
For the two positive contexts above, \(C\) may produce \(a\) in the
first and \(E\) in the second, with \(A\) irrelevant to each.
Difference can still identify a causal contribution in the particular
matched comparison without asserting that no other route produces
the same phenomenon elsewhere
\parencite[III.X.1--3]{Mill1882SystemLogicPG}.

The following finite constructions make two boundaries exact.
They are editorial illustrations, not historical axiomatizations of
either author's whole method.

\begin{hbtheorem}{thm:p15-wm-causal-ambiguity}{proved=annex:p15.wm.causal-ambiguity}{Agreement and covariation can leave the cause unsettled}
Let the complete context consist of some subset of the five distinct
circumstances \(A,B,C,D,E\), and let \(a\) occur exactly when \(C\)
or \(E\) is present. The positive contexts \(\{A,B,C\}\) and
\(\{A,D,E\}\) share exactly \(A\), but removing \(A\) from either
does not remove \(a\).

For a distinct binary example let \(U\) take values zero and one,
each with probability \(1/2\). Consider the assignments
\[
 \text{direct model: }X=U,\quad Y=X;
 \qquad
 \text{common-source model: }X=U,\quad Y=U.
\]
The observational law of \((X,Y)\) is the same in both models.
Replacing the equation for \(X\) by the assignment \(X=0\), while
leaving the other equations and the law of \(U\) fixed, gives
different laws of \(Y\): always zero in the direct model, and a
fair binary outcome in the common-source model.
\end{hbtheorem}

In the first construction plurality, not a missing common circumstance,
defeats the overstrong inference. In the second, the stipulated equation
replacement is a modern intervention interpretation. Observed concordance
does not justify that interpretation or decide between its two models.
Neither construction proves that a particular scientific experiment
satisfies the isolation conditions of Mill's method.

\begin{hbtheorem}{thm:p15-wm-finite-continuation}{proved=annex:p15.wm.finite-continuation}{Exact finite fit does not determine extrapolation}
Let \(m\geq1\), let \(x_1,\ldots,x_m\) be distinct real numbers, and
let \(y_1,\ldots,y_m\) be arbitrary real numbers. There is a polynomial
\[
 p(x)=\sum_{i=1}^m y_i
             \prod_{\substack{1\leq j\leq m\\j\ne i}}
                    \frac{x-x_j}{x_i-x_j}
\]
which fits every pair. For each real \(c\), the polynomial
\[
 p_c(x)=p(x)+c\prod_{i=1}^m(x-x_i)
\]
fits the same pairs. At any real \(t\) different from every \(x_i\),
the value \(p_c(t)\) can be any prescribed real number as \(c\) varies.
Empty products equal one.
\end{hbtheorem}

For an explicit exact record, the points \((0,0)\) and \((1,1)\)
are fitted both by \(p_0(x)=x\) and by \(p_1(x)=x+x(x-1)\).
They predict two and four, respectively, at \(x=2\).
Choosing a class, degree restriction, error tolerance or physical mechanism
can exclude alternatives, but those are additional constraints.
This example supports the distinction between fit and extension; it does
not prove that every inductive conclusion is equally warranted or that
consilience supplies no evidence.

The neighboring Bayesian and learning chapters specify some such additional
structures explicitly: observation models and priors in
\cref{sec:p15-bayesian-updating-prediction}, or hypothesis classes and
sampling assumptions in \cref{sec:p15-occam-learning}.
They are not translations which settle the Whewell--Mill dispute by
renaming colligation a posterior distribution or consilience an
independent sample.

The neighboring Hume treatment in
\cref{sec:p15-hume-and-the-problem-of-induction} separates a philosophical
account of the passage beyond experience from a conditional model's
consequences. The distinction matters here too: organizing a hypothesis,
isolating a causal contribution and fitting a predictive law are different
tasks, not three demonstrations that no unobserved exception is possible.

\begin{proofboundary}{def:p15-wm-scope}{Historical arguments and the formal illustrations}
The two finite results prove only their declared causal assignments and
polynomial claims. Whewell's historical truth tests and Mill's experimental
canons retain their attributed philosophical assumptions and qualifications.
This primary-text comparison does not supply the whole books, a modern
physical reconstruction of all their scientific examples, or the distinct
complete 2011 historical treatment by Malcolm Forster
\parencite{Forster2011WhewellMill}.
\end{proofboundary}

\index{Whewell, William}\index{Mill, John Stuart}
\index{induction!scientific}\index{colligation}\index{consilience}
\index{causal inference!Mill's methods}\index{extrapolation!finite-fit ambiguity}
